% Standalone manuscript. Compile twice with pdflatex, or once with tectonic.
% Set \notebooknotesfalse below to omit the two documentary appendices.
\documentclass[11pt,letterpaper]{article}
\usepackage[margin=1in]{geometry}
\usepackage{amsmath,amssymb,booktabs}
\usepackage[hidelinks]{hyperref}
\hypersetup{pdftitle={NOMAD: A Taylor solver for multivariate polynomial differential equations},
  pdfauthor={John Atkinson},pdfsubject={Notebook derivation and implementation notes}}
\newif\ifnotebooknotes
\notebooknotestrue
\newcommand{\RR}{\mathbb{R}}
\newcommand{\NN}{\mathbb{N}_0}
\newcommand{\dd}{\mathrm{d}}
\newcommand{\ii}{\mathrm{i}}
\newcommand{\D}{\mathcal{D}_F}
\newcommand{\src}[1]{\par\smallskip\noindent{\small\itshape Notebook, #1.}\par\medskip}
\newcommand{\leaf}[2]{\par\smallskip\begingroup\interlinepenalty=10000\noindent\textbf{#1}\enspace #2\par\endgroup}
\newcommand{\notegroup}[1]{\subsection*{#1}\addcontentsline{toc}{subsection}{#1}}
\newcommand{\file}[1]{\texttt{\detokenize{#1}}}
\DeclareMathOperator{\sech}{sech}
\DeclareMathOperator{\csch}{csch}
\setcounter{tocdepth}{2}
\setlength{\emergencystretch}{1.3em}
\setlength{\parindent}{1.2em}
\setlength{\parskip}{2pt}
\raggedbottom
\begin{document}
\begin{center}
{\LARGE\bfseries NOMAD\par}
\vspace{.4em}
{\Large A Taylor solver for multivariate polynomial\\[.15em] differential equations\par}
\vspace{.8em}
{\large John Atkinson\par}
\vspace{.3em}
{\small Whitepaper draft from the 2016--2017 notebooks\par}
\end{center}
\begin{abstract}
For an autonomous polynomial system, every time derivative of the solution is
itself a polynomial in the current state. NOMAD constructs those polynomials
before a time step is taken. It stores the exponent of each monomial in one
column and the coefficients for all components of the system in a matching
column. A single differentiation combines an old column with a column of the
vector field, subtracts one from the exponent of the differentiated variable,
and adds contributions with equal exponents. This paper derives both the explicit
formula for the \(n\)th derivative and the sliding matrix procedure for computing
it. A two-variable example follows the columns through to a Taylor step.
The final part of the notebook gives an exponential-quotient representation for
circular and hyperbolic functions of polynomial arguments. Finding a comparable
higher-derivative algorithm for that representation remained open in the notes.
\end{abstract}

\section{The problem in the notebook}\label{sec:problem}
The starting point was a question about an ordinary multivariate polynomial.
If a table gives the exponent of each variable in each term, what changes in
that table when the polynomial is differentiated? The notebook first answers
the question for a gradient and a Hessian. It then uses a polynomial for each
right-hand side of a system
\begin{equation}\label{eq:ivp}
   \dot x=F(x),\qquad x(t_0)=z,\qquad x=(x_1,\ldots,x_M)^\mathsf{T}\in\RR^M .
\end{equation}
Now differentiating \(x_j(t)\) inserts another polynomial, \(F_j\). The working
problem on p.~11 is to find \(\dd^n x_\ell/\dd t^n\) directly from the
coefficients and exponents of those first derivatives. Once constructed, this
expression can be evaluated at any initial state for the same field.

The cover of the photographic notebook~\cite{notebook} calls NOMAD the “N-th Order M-Dimensional Nonlinear Diff Eq Solver.”
The derivation began during 2016--2017, and the clean review in the notebook is
dated 27 August 2017. We use \(N\) for a bound on each variable's exponent in
\(F\), \(n\) for a time derivative, and \(p\) for the order retained in a
numerical step. The equations in \eqref{eq:ivp} are first order in time.

\paragraph{Relation to Taylor integration.}
The operator \(F\cdot\nabla\) that appears below is the Lie derivative along
the field. Iterating it on coordinate functions gives the classical Lie-series
description of the flow \cite{pal}. High-order Taylor integrators can also
compute series coefficients at a given state by automatic differentiation
\cite{jorba}. The object constructed in this notebook is a table of
polynomials \emph{in the state}, one for every requested order. Its
construction cost can therefore be shared across time steps and initial
conditions. The question addressed here is how to build and collect those
tables from the monomials of \(F\).
\src{pp.~2--11}

\section{The coefficient and exponent tables}\label{sec:encoding}
For a column \(k\in\NN^M\), write
\(x^k=\prod_{\beta=1}^M x_\beta^{k_\beta}\), with a factor of exponent zero
equal to one, including at \(x_\beta=0\). Let \(e_j\) be the \(j\)th unit
vector and \(\delta_{\beta j}=(e_j)_\beta\). If the largest individual power
is \(N\), the full table \(K\) contains \(T=(N+1)^M\) columns \(k_i\).
The last coordinate varies fastest, and each digit runs from \(N\) down to
zero. Equivalently, for \(1\leq i\leq T\),
\begin{equation}\label{eq:digit}
  K_{\beta i}=N-
  \left(\left\lfloor\frac{i-1}{(N+1)^{M-\beta}}\right\rfloor
        \bmod (N+1)\right).
\end{equation}
In row \(\beta\), each value is repeated \((N+1)^{M-\beta}\) times per
run, and the full descending run is repeated \((N+1)^{\beta-1}\) times.
This typesets the repeated blocks of ones on pp.~3 and 10. The corner
\(x_1^N\cdots x_M^N\) belongs to the table and has total degree \(MN\).

For a scalar polynomial \(P(x)=\sum_i a_i x^{k_i}\), the gradient and Hessian
on p.~2 follow by retaining the original column index in each sum,
\begin{align}
 \partial_w P
 &=\sum_{i:K_{wi}>0} a_iK_{wi}\,x^{k_i-e_w},\label{eq:grad}\\
 \partial_w^2 P
 &=\sum_{i:K_{wi}\ge2}a_iK_{wi}(K_{wi}-1)\,x^{k_i-2e_w},\label{eq:hpure}\\
 \partial_y\partial_w P
 &=\sum_{\substack{i:K_{wi}>0\\K_{yi}>0}}
    a_iK_{wi}K_{yi}\,x^{k_i-e_w-e_y},\qquad y\ne w .\label{eq:hmixed}
\end{align}
The notebook calls the retained coefficients \(b_{wi}\) and \(c_{ywi}\).
The column counts before coefficient zeros are taken into account are
\(N(N+1)^{M-1}\) for \eqref{eq:grad},
\(\max(N-1,0)(N+1)^{M-1}\) for \eqref{eq:hpure}, and
\(N^2(N+1)^{M-2}\) for \eqref{eq:hmixed} when \(M\ge2\).
Stacking the first derivatives gives \(\nabla P\), and arranging the second
derivatives by \(y,w\) gives the Hessian.

With a common exponent table and a real \(M\times T\) coefficient matrix \(A\), the field has components
\begin{equation}\label{eq:field}
 F_\ell(x)=\sum_{i=1}^{T} A_{\ell i}\,x^{k_i}.
\end{equation}
A zero entry excludes a monomial in one equation. A whole column may be
discarded when \(A_{:i}=0\). Below \(S\) denotes the number of columns kept,
relabeled \(1,\ldots,S\), with \(A,K\) denoting this reduced pair. The same
formulas hold for the full pair. The all-constant term has \(k_i=0\) and may
have a nonzero coefficient.

\paragraph{The repeated state array.}
The working notes write \(X^\beta_\alpha=x_\beta\) for every column
\(\alpha\). Raising \(X\) to \(K\) elementwise and taking a product down each
column gives the vector of monomials \(m_K(x)=(x^{k_1},\ldots,x^{k_S})^\mathsf{T}\).
Thus \(F(x)=A\,m_K(x)\). The symbol \(\dd x^\ell_\alpha/\dd t\) on pp.~25--28
labels the contribution \(A_{\ell\alpha}x^{k_\alpha}\) to \(\dot x_\ell\).
All columns are evaluated at the same state. This convention becomes essential
when the notes regroup sums over \(\alpha\).
\src{pp.~2--4, 10--11, 25--29, 39, 45--46}

\section{Repeated time differentiation}\label{sec:derivatives}
For any polynomial \(G\), the chain rule along \eqref{eq:ivp} gives
\[
  \frac{\dd}{\dd t}G(x(t))=(\D G)(x(t)),\qquad
  \D G=\sum_{j=1}^M F_j\,\partial_jG .
\]
In the handwritten derivation this identity emerges by peeling off the factors
of \(x^u=x_1^{u_1}\cdots x_M^{u_M}\) one at a time. Differentiating the
\(j\)th factor pulls down \(u_j\) and leaves the other factors intact. Inserting
\eqref{eq:field} gives
\begin{equation}\label{eq:one}
 \D(x^u)
 =\sum_{j:u_j>0}u_j x^{u-e_j}F_j(x)
 =\sum_{j:u_j>0}\sum_{i=1}^S u_j A_{ji}\,x^{u+k_i-e_j}.
\end{equation}
When \(u_j=0\), this branch ends before forming an exponent. This also
defines the computation on coordinate planes. The extra undifferentiated
product in the trial on p.~12 disappears in the correction on pp.~13--14.

\subsection{From the second derivative to the fourth}
Choose the initial column \(\alpha\). Each further differentiation chooses a
coordinate \(j_r\) and a refreshed field column \(i_r\). Write
\[
 u_0=k_\alpha,\qquad
 u_r=k_\alpha+\sum_{s=1}^{r}(k_{i_s}-e_{j_s}).
\]
In the following displays, \(\sum^{*}\) omits a choice whenever
\((u_{r-1})_{j_r}=0\). A coordinate index runs from 1 to \(M\), and a column
index from 1 to \(S\). With \(x=x(t)\), the first four derivatives are
\begin{align}
 x_\ell' &= \sum_\alpha A_{\ell\alpha}x^{u_0},\nonumber\\
 x_\ell'' &= \sum_{\alpha,j_1,i_1}^{*}
  A_{\ell\alpha}\,K_{j_1\alpha} A_{j_1 i_1}\,x^{u_1},\label{eq:d2}\\
 x_\ell''' &= \sum_{\alpha,j_1,i_1,j_2,i_2}^{*}
  A_{\ell\alpha}\,K_{j_1\alpha} A_{j_1 i_1}
    (K_{j_2\alpha}+K_{j_2i_1}-\delta_{j_2 j_1}) A_{j_2i_2}\,x^{u_2},
    \label{eq:d3}\\
 x_\ell^{(4)} &= \sum_{\alpha,j_1,i_1,j_2,i_2,j_3,i_3}^{*}
 A_{\ell\alpha}\prod_{r=1}^{3}\big[(u_{r-1})_{j_r} A_{j_ri_r}\big]\,
 x^{u_3}. \label{eq:d4}
\end{align}
The new multiplier in \eqref{eq:d4} reads
\(K_{j_3\alpha}+K_{j_3i_1}+K_{j_3i_2}-\delta_{j_3j_1}-\delta_{j_3j_2}\).
On p.~19 the fourth derivative is a prediction. Pages 20--21 verify it by
naming the entire old exponent \(\hat k\) and applying the product rule to
\(x^{\hat k}\). This calculation identifies the exponent one must carry
forward at \emph{every} differentiation.
\src{pp.~4--6, 12--21, 23--24}

\subsection{The explicit formula for any derivative}
The notebook abbreviates the \(2n-1\) nested sums using \(J(n)\) and \(I(n)\).
The coordinate and column choices range over the ordered tuples
\[
  (j_1,\ldots,j_{n-1})\in\{1,\ldots,M\}^{n-1},\qquad
  (i_1,\ldots,i_{n-1})\in\{1,\ldots,S\}^{n-1}.
\]
An empty choice contributes once and an empty product equals one. The time
derivative of order \(n\ge1\) is
\begin{equation}\label{eq:path}
 \boxed{\displaystyle
 x_\ell^{(n)}(t)=\sum_{\alpha,J(n),I(n)}^{*}
 A_{\ell\alpha}
 \left(\prod_{r=1}^{n-1}(u_{r-1})_{j_r} A_{j_ri_r}\right)
 x(t)^{u_{n-1}} .}
\end{equation}
To recover the scalar-index expression boxed on p.~30, set \(i_0=\alpha\),
\(j_0=\ell\), and \(U(s-1)=0\) at \(s=0\), \(1\) at integer \(s\ge1\).
Define
\[
  B_\beta(r)=\sum_{s=0}^{r}
       \big(K_{\beta i_s}-\delta_{\beta j_s} U(s-1)\big).
\]
Each summand of \eqref{eq:path} is the notebook's three factors,
\begin{equation}\label{eq:three}
 \underbrace{\prod_{r=0}^{n-1} A_{j_r i_r}}_{\text{field coefficients}}
 \quad\underbrace{\prod_{r=1}^{n-1} B_{j_r}(r-1)}_{\text{differentiated powers}}
 \quad\underbrace{\prod_{\beta=1}^M x_\beta^{B_\beta(n-1)}}_{\text{remaining monomial}}.
\end{equation}
The step function suppresses the Kronecker subtraction at the initial column.
For \(n=1\), \eqref{eq:path} therefore reduces to the input field.

\paragraph{Proof.}
The case \(n=1\) is \eqref{eq:field}. In any term at order \(n\) the
coefficient is independent of time. Applying \eqref{eq:one} to the remaining
monomial extends its choices by \((j_n,i_n)\), multiplies its coefficient by
\((u_{n-1})_{j_n}A_{j_n i_n}\), and leaves exponent \(u_n\). These are exactly
the terms and factors of \eqref{eq:path} at order \(n+1\). This proves the
formula by induction. Before zero factors are considered there are
\(S(MS)^{n-1}\) choices for a nonempty field. Many lead to the same monomial.
\src{pp.~5--7, 19--30, 54}

\section{Collecting derivatives into matrices}\label{sec:recurrence}
The move from the expanded expression to the algorithm occurs on pp.~31--33.
Performing the innermost column sum in the second derivative gives
\[
 x_\ell''=\sum_{\alpha,j:K_{j\alpha}>0}
 \underbrace{A_{\ell\alpha}K_{j\alpha}x^{k_\alpha-e_j}}
              _{\partial_j(A_{\ell\alpha}x^{k_\alpha})}
 \underbrace{\sum_i A_{ji}x^{k_i}}_{F_j(x)}.
\]
For the third derivative the same factorization works with each term of
\(x_\ell''\). The inner sums over \(j_1,i_1\) have already determined a
coefficient and a column \(k_\alpha+k_{i_1}-e_{j_1}\). Those two objects
contain everything needed for the next differentiation. In particular,
terms with equal exponents can be added \emph{before} differentiating again.

Let \(H_n=\D^{\,n-1}F\), so \(H_n(x(t))=x^{(n)}(t)\). With each exponent \(q_\alpha\) in
\(\widehat K^{(n)}\) store its vector of coefficients
\(c_\alpha=\widehat A^{(n)}_{:\alpha}\). Then
\begin{equation}\label{eq:table}
  H_1=F,\qquad H_n(x)=\sum_{\alpha=1}^{S_n}c_\alpha x^{q_\alpha},
  \qquad H_{n+1}=\D H_n .
\end{equation}
Initialize \((\widehat A^{(1)},\widehat K^{(1)})=(A,K)\). For each choice of \((\alpha,j,i)\) with \((q_\alpha)_j>0\),
\eqref{eq:one} gives the update
\begin{equation}\label{eq:rec}
 \boxed{\displaystyle
 v=q_\alpha+k_i-e_j,\qquad
 \Delta c=c_\alpha\,(q_\alpha)_j\,A_{ji}.}
\end{equation}
For each distinct \(v\), its new column is
\begin{equation}\label{eq:group}
 \widehat A^{(n+1)}_{:v}
 =\sum_{\substack{\alpha,j,i:\ (q_\alpha)_j>0\\ q_\alpha+k_i-e_j=v}}
     c_\alpha (q_\alpha)_j A_{ji}.
\end{equation}
Here \(v\) is the exponent used to look up the column. In the review on
p.~8, the multiplying exponent is written with row \(\ell\). The raw
derivation on pp.~53 and 55 and the product rule both give row \(j\).
The output component \(\ell\) selects an entry of \(c_\alpha\).
\src{pp.~8--9, 26--27, 31--37, 44, 53--55}

\subsection{Sliding the original table across the current table}\label{sec:slide}
The notebook draws a circled plus between \(\widehat K^{(n)}\) and the shifted
matrix \(K-e_j{\bf1}^{\mathsf T}\). For a fixed \(j\), the first column of the
shifted matrix first touches the last column of the current matrix.
Both arrays keep their orientation as the shifted matrix moves across. Each
overlap gives a list of pairs \((\alpha,i)\). For two tables of four columns,
the entire handwritten index list on p.~48 reads as follows.
\begin{center}
\begin{tabular}{@{}rll@{}}
\toprule
Offset \(\alpha-i\)& Current columns \(\alpha\)& Field columns \(i\)\\
\midrule
 3&4&1\\
 2&3, 4&1, 2\\
 1&2, 3, 4&1, 2, 3\\
 0&1, 2, 3, 4&1, 2, 3, 4\\
 \(-1\)&1, 2, 3&2, 3, 4\\
 \(-2\)&1, 2&3, 4\\
 \(-3\)&1&4\\
\bottomrule
\end{tabular}
\end{center}
For \(P=S_n\) current columns and \(Q=S\) field columns, both positive, take
\(d=P-1,P-2,\ldots,1-Q\). At offset \(d\), take
\[
 \max(1,1-d)\le i\le\min(Q,P-d),\qquad \alpha=i+d.
\]
Each pair has exactly one offset. There are \(P+Q-1\) alignments and \(PQ\)
pairs per coordinate. With \(m=\min(P,Q)\) and \(L=\max(P,Q)\), the notebook's
count is \(2\sum_{r=1}^{m-1}r+m(L-m+1)=PQ\). For \(P=6,Q=4\) the overlap
lengths are \(1,2,3,4,4,4,3,2,1\), a total of 24 pairs.

The analogy with convolution has a precise interpretation. Index the old
coefficient vector by its exponent \(u\), writing it as \(c(u)\), and index
the \(j\)th input coefficient by its exponent \(w\), writing \(a_j(w)\).
Use zero for an absent coefficient. The coefficient of \(x^v\) after one
differentiation is
\begin{equation}\label{eq:conv}
 \sum_{j=1}^M\;\sum_{\substack{b,w\in\NN^M\\b+w=v}}
       (b_j+1)c(b+e_j)a_j(w).
\end{equation}
For each \(j\), polynomial multiplication convolves the exponent-indexed
coefficients of \(\partial_j H_n\) and \(F_j\). The drawn slide enumerates
pairs of \emph{stored columns}. A slide offset can contain several different
exponents, and the same exponent can occur at several offsets or coordinates.
Equation~\eqref{eq:group} collects by the full exponent vector across all of them.

\paragraph{Connection with the Carleman matrix.}
The same coefficients can be indexed by the monomial being differentiated
and the monomial it produces. Write
\[
 L_{uv}=\sum_{j:u_j>0}u_j\,a_j(v-u+e_j),\qquad
 \D(x^u)=\sum_v L_{uv}x^v ,
\]
with \(a_j(w)=0\) for any missing exponent. For fields with \(F(0)=0\)
this entry is the monomial Carleman matrix formula in
Chermnykh \cite[Prop.~4, Eq.~(5)]{carleman}. A constant field term in the
present problem is included by allowing \(w=0\). In this notation NOMAD's
grouped update is \(c_{\rm new}(v)=\sum_u c(u)L_{uv}\). It constructs
the entries reached from the current derivative and truncates the solution
in time order \(p\). This gives a point of comparison for the notebook's
sliding construction.

\paragraph{One derivative in the notebook's order.}
The following is the algorithm of p.~53, written for current width \(P\) and
original width \(Q\). All coefficient additions are vector additions.
\begin{enumerate}
\item Start an empty table indexed by \(v\), with missing coefficients equal to
the zero vector. For each \(j=1,\ldots,M\), traverse the offsets above.
\item For each aligned pair \((\alpha,i)\), read \((q_\alpha)_j\) and
\(A_{ji}\). If either is zero, pass to the next pair. Otherwise form \(v\) and
add \(\Delta c\) from \eqref{eq:rec} to the entry for \(v\).
\item Remove entries whose accumulated vector is zero. Sort the remaining
exponents in descending lexicographic order, moving each coefficient vector
with its exponent. These are \(\widehat K^{(n+1)},\widehat A^{(n+1)}\).
\end{enumerate}
An empty resulting table represents the identically zero derivative. Its later
derivatives vanish in exact arithmetic. If the original table is empty, the
solution is the constant state \(z\). A zero found by evaluating a table at a particular state permits no such inference.

\paragraph{How large can a table become?}
Pages 7 and 34 estimate a dense allocation of \((2^{n-1}N+1)^M\) slots
by successively doubling a bound on each exponent. This is a loose upper
bound. Each coordinate of \(u_{n-1}\) in \eqref{eq:path} is at most \(nN\), which gives at most \((nN+1)^M\)
distinct exponents. If \(D\ge1\) bounds the \emph{total} degree of \(F\),
each surviving differentiation adds at most \(D-1\) to total degree, so
\(\deg H_n\le 1+n(D-1)\). The actual table has only the \(S_n\) keys that
survive collection. Building it requires at most \(MS_{n-1}S\) candidate
vectors for \(n\ge2\), before the zero tests. This distinction between candidate
products and surviving monomials motivated the slide.
\src{pp.~7--9, 33--53}

\section{A hand computation in two variables}\label{sec:example}
The example that connects the early algebra to the implementation is
\begin{equation}\label{eq:xy}
 \dot x=y^2-x,\qquad \dot y=x^2-y .
\end{equation}
On p.~22 the full exponent matrix and its selected coefficients are
\[
 K_{\rm full}=
 \begin{pmatrix}2&2&2&1&1&1&0&0&0\\2&1&0&2&1&0&2&1&0\end{pmatrix},
 \qquad
 A_{\rm full}=
 \begin{pmatrix}0&0&0&0&0&-1&1&0&0\\0&0&1&0&0&0&0&-1&0\end{pmatrix}.
\]
Keeping columns \(3,6,7,8\) gives
\begin{equation}\label{eq:xy1}
 K=\begin{pmatrix}2&1&0&0\\0&0&2&1\end{pmatrix},\qquad
 A=\begin{pmatrix}0&-1&1&0\\1&0&0&-1\end{pmatrix}.
\end{equation}
For \(j=1\) the notebook slides the four columns
\((1,0),(0,0),(-1,2),(-1,1)\) across \(K\). A shifted field column can
contain \(-1\) before the current exponent is added. At offset \(d=-1\) the shifted table extends one column to the left. The
three overlapping pairs are aligned below.
\[
 \begin{array}{r@{\quad}ccccc}
  \text{current } q_\alpha & & (2,0) & (1,0) & (0,2) & (0,1)\\
  \text{shifted } k_i-e_1  & (1,0) & (0,0) & (-1,2) & (-1,1) & \\
  \hline
  (\alpha,i) & & (1,2) & (2,3) & (3,4) &\\
  q_\alpha+k_i-e_1 & & (2,0) & (0,2) & (-1,3) &
 \end{array}
\]
The third pair has \((q_3)_1=0\) and is discarded. The second pair
\(\alpha=2,\ i=3,\ j=1\) gives \(q_\alpha=(1,0)\) and
\(k_i-e_1=(-1,2)\). Their sum is \((0,2)\), with vector contribution
\((-1,0)^\mathsf{T}\). The coordinate \(j=2\) contributes
\((-2,0)^\mathsf{T}\) to this same key from \(\alpha=3,\ i=4\). This is
the \(-3y^2\) in \(x''\). Comparing or rejecting the shifted field column
before adding \(q_\alpha\) would lose the first contribution.

Here are the two passes in polynomial form, with one entry per output component.
\begin{center}
\begin{tabular}{@{}lll@{}}
\toprule
Differentiated coordinate & Contribution to \(x''\) & Contribution to \(y''\)\\
\midrule
\(x\) & \(x-y^2\) & \(2xy^2-2x^2\)\\
\(y\) & \(2x^2y-2y^2\) & \(y-x^2\)\\
\bottomrule
\end{tabular}
\end{center}
After collecting and sorting, pp.~50--51 give
\begin{equation}\label{eq:xy2}
 \widehat K^{(2)}=\begin{pmatrix}2&2&1&1&0&0\\1&0&2&0&2&1\end{pmatrix},
 \qquad
 \widehat A^{(2)}=\begin{pmatrix}2&0&0&1&-3&0\\0&-3&2&0&0&1\end{pmatrix}.
\end{equation}
Reading off the columns gives
\begin{equation}\label{eq:xyD2}
 x''=2x^2y+x-3y^2,\qquad y''=2xy^2+y-3x^2 .
\end{equation}
Direct differentiation gives the same expressions from
\(x''=2y(x^2-y)-(y^2-x)\) and \(y''=2x(y^2-x)-(x^2-y)\).
The six columns of \eqref{eq:xy2} become the current table for the next
slide. The four columns in \eqref{eq:xy1} remain the field table.
\src{pp.~22, 45, 47--53}

\section{From derivative tables to a time step}\label{sec:taylor}
With \(z=x(t_0)\), the Taylor series on p.~7 reads
\begin{equation}\label{eq:series}
 x(t_0+h)=z+\sum_{n=1}^{\infty}\frac{h^n}{n!}H_n(z).
\end{equation}
A polynomial field is analytic, so the solution has such a series locally.
Its convergence radius depends on the system and initial state. For instance,
\(\dot x=x^2\) gives \(H_n(z)=n!\,z^{n+1}\) for \(n\ge1\). The series is
geometric and sums to \(z/(1-hz)\) for \(|hz|<1\) when \(z\ne0\). The finite
formula for each derivative and the convergence interval of their infinite sum
are separate conclusions.

\subsection{Distributing the factorial}
Pages 54--57 describe a practical difficulty. Raw derivative coefficients and
\(n!\) can both be large well before their quotient is needed. Let
\(C_n(x)=H_n(x)/n!\). Linearity of \(\D\) gives
\begin{equation}\label{eq:scale}
 C_1=F,\qquad C_{n+1}=\frac{\D C_n}{n+1}.
\end{equation}
Thus, in passing from order \(n\) to \(n+1\), the scaled version of
\eqref{eq:rec} adds \(c_\alpha(q_\alpha)_j A_{ji}/(n+1)\), where
\(c_\alpha\) is already a coefficient of \(C_n\). Only the old derivative
coefficients are scaled. The original field \(A\) is used in every pass.
Page 57 instead stores a partially scaled table representing
\(H_n/(n-1)!\). Its series term still requires a final division by \(n\).
Equation~\eqref{eq:scale} incorporates that last division in the stored table.
The divisions \(1/2,1/3,\ldots,1/n\) accumulate to \(1/n!\).
Scaling controls one source of growth. Rounding, cancellation and the size of
the remaining coefficients still depend on the field and the order.

\subsection{Evaluation and restart}
Precompute \(C_1,\ldots,C_p\) for the chosen field and retain their exponent
tables. With state \(z_m\) at \(t_m\) and step \(h_m=t_{m+1}-t_m\), set
\begin{equation}\label{eq:step}
 \boxed{\displaystyle z_{m+1}=z_m+\sum_{n=1}^{p}h_m^n C_n(z_m).}
\end{equation}
The expansion state is fixed throughout the sum. Only after completing it is
the series restarted at \(z_{m+1}\). For an exact starting state and bounded
\((p+1)\)st derivative over the step, the local error is \(O(|h_m|^{p+1})\)
by Taylor's theorem. On a fixed interval where the solution and steps stay in
a compact neighborhood, smoothness gives the usual global
\(O((\max_m |h_m|)^p)\) bound in exact arithmetic for sufficiently small steps.
The adaptive feature on p.~34 concerns the number of stored monomials. The
notebook's series takes the step size as an input.

As an arithmetic check added for this edition, take \((x,y)=(1,2)\) in
\eqref{eq:xy}. The first four derivatives at that point are
\[
 \begin{array}{c|rrrr}
    n&1&2&3&4\\ \hline
    x^{(n)}&3&-7&37&-91\\
    y^{(n)}&-1&7&-3&-49
 \end{array}.
\]
For example \(x'''=2(y')^2+2yy''-x''=37\).
The second-order step is
\((1+3h-\tfrac72h^2,\ 2-h+\tfrac72h^2)\).
An independent check uses coefficient matching. If
\(x(t_0+h)=\sum u_n h^n\) and \(y(t_0+h)=\sum v_n h^n\), then
\[
 (n+1)u_{n+1}=\sum_{k=0}^n v_kv_{n-k}-u_n,\qquad
 (n+1)v_{n+1}=\sum_{k=0}^n u_ku_{n-k}-v_n.
\]
Starting from \((u_0,v_0)=(1,2)\), these rational recurrences agree with
\(H_n(1,2)/n!\) through order eight. They offer a check of the factorials
independent of the sliding indices.
\src{pp.~7--9, 50--57}

\clearpage
\section{Polynomial arguments inside transcendental terms}\label{sec:trans}
The last eight photographs change the question. How much of the structure
of the input field can be retained if a monomial is multiplied by a circular
or hyperbolic function? Page 58 begins with the identities
\begin{align*}
 \sin z&=(e^{\ii z}-e^{-\ii z})/(2\ii),  &
 \sinh z&=(e^z-e^{-z})/2=-\ii\sin(\ii z),\\
 \cos z&=(e^{\ii z}+e^{-\ii z})/2, &
 \cosh z&=(e^z+e^{-z})/2=\cos(\ii z),\\
 \tan z&=\frac{e^{\ii z}-e^{-\ii z}}{\ii(e^{\ii z}+e^{-\ii z})}, &
 \tanh z&=\frac{e^z-e^{-z}}{e^z+e^{-z}}=-\ii\tan(\ii z).
\end{align*}
The even and odd exponential pairs provide numerators and denominators.
Choosing the signs and the arguments selects the desired function.

\subsection{The expression on the last page}
The shorthand introduced on pp.~60--61 is itself the original scalar
polynomial. With a coefficient row \(u\) and exponent table \(V\), write
\begin{equation}\label{eq:Y}
 Y(x;u,V)=\sum_{\gamma=1}^{|u|}u_\gamma
                  \prod_{\beta=1}^M x_\beta^{V_{\beta\gamma}}.
\end{equation}
Here \(|u|\) counts the columns. The manuscript abbreviates this as
\(Y(u,V)\). In one term use four such arguments,
\[
 C=Y(x;c,P),\quad D=Y(x;d,P),\quad
 B=Y(x;f,Q),\quad G=Y(x;g,Q).
\]
The two numerator polynomials share an exponent table \(P\), and the two
denominator polynomials share \(Q\). In dense form their respective column
counts are \((N_P+1)^M\) and \((N_Q+1)^M\), where the two bounds on the
individual powers may differ. Complex coefficients are allowed in these four arguments.

For each component \(\ell\) and monomial \(\alpha\), choose selectors
\(r_{\ell\alpha},s_{\ell\alpha}\in\{0,1\}\) and define
\begin{equation}\label{eq:ratio}
 R_{\ell\alpha}(x)=
 \frac{\exp C+(-1)^{r_{\ell\alpha}}\exp D}
      {\exp B+(-1)^{s_{\ell\alpha}}\exp G}.
\end{equation}
Page 65 writes the inner coefficient rows with a component index and leaves
their dependence on the outer term implicit. The display allows them to be
selected per term. Reusing the same rows recovers the shared-argument case.
Page 59 first pairs \(C,D\) as \(Y_1,-Y_1\) and \(B,G\) as \(Y_2,-Y_2\).
The four-polynomial form on p.~65 keeps the pairs available separately. In
this notation the general equation on that page reads
\begin{equation}\label{eq:general}
 \boxed{\displaystyle
 \dot x_\ell=\sum_{\alpha=1}^{(N+1)^M}
 \big[a_{\ell\alpha}+b_{\ell\alpha}R_{\ell\alpha}(x)^{h_{\ell\alpha}}\big]\,
 x^{k_\alpha}.}
\end{equation}
The power \(h_{\ell\alpha}\) acts on the quotient and \(b_{\ell\alpha}\)
multiplies the result. Setting \(b_{\ell\alpha}=0\) omits the transcendental
part of that term, so its quotient need never be evaluated. Setting
\(a_{\ell\alpha}=0\) isolates the transcendental term times its monomial,
as in the lower box on p.~65. For an active quotient the denominator must
be nonzero. The notebook leaves the range of \(h_{\ell\alpha}\) unspecified.
All the named functions below use \(h_{\ell\alpha}=1\). Other choices require
a domain on which the chosen power is defined.

\subsection{The twelve function selections}
Let \(\theta=Y(x;u,V)\) have real coefficients, let \(w\) be a real
amplitude, and put \(E_0(z)=e^z+e^{-z}\), \(E_1(z)=e^z-e^{-z}\).
Take \(C=v,\ D=-v,\ B=z,\ G=-z\) in \eqref{eq:ratio}. Table~\ref{tab:selectors}
gives choices for which \(bE_r(v)/E_s(z)\) is the indicated function times
\(w\). Choosing zero arguments with the even selector gives \(E_0(0)=2\).
The factors of \(\ii\) in \(b\) distinguish a sine from its reciprocal.

\begin{table}[!ht]
\centering
\caption{Sign and amplitude choices in the exponential quotient for real \(\theta\).}
\label{tab:selectors}\smallskip
\begin{tabular}{@{}lccccc@{}}
\toprule
Function of \(\theta\)&\(v\)&\(z\)&\(r\)&\(s\)&\(b\)\\
\midrule
\(\sin\)  &\(\ii\theta\)&0&1&0&\(-\ii w\)\\
\(\sinh\) &\(\theta\)&0&1&0&\(w\)\\
\(\cos\)  &\(\ii\theta\)&0&0&0&\(w\)\\
\(\cosh\) &\(\theta\)&0&0&0&\(w\)\\
\(\sec\)  &0&\(\ii\theta\)&0&0&\(w\)\\
\(\sech\) &0&\(\theta\)&0&0&\(w\)\\
\(\csc\)  &0&\(\ii\theta\)&0&1&\(\ii w\)\\
\(\csch\) &0&\(\theta\)&0&1&\(w\)\\
\(\tan\)  &\(\ii\theta\)&\(\ii\theta\)&1&0&\(-\ii w\)\\
\(\tanh\) &\(\theta\)&\(\theta\)&1&0&\(w\)\\
\(\cot\)  &\(\ii\theta\)&\(\ii\theta\)&0&1&\(\ii w\)\\
\(\coth\) &\(\theta\)&\(\theta\)&0&1&\(w\)\\
\bottomrule
\end{tabular}
\end{table}

For a sine, the numerator arguments have opposite, purely imaginary
coefficient rows, while \(f=g=0\) and \(s=0\). This makes the term
\(-\ii w(e^{\ii\theta}-e^{-\ii\theta})/2=w\sin\theta\).
For a hyperbolic cosecant, the numerator equals two, and opposite real
denominator arguments give \(2w/(e^\theta-e^{-\theta})=w\csch\theta\).
Page 63 writes equal real rows in the latter denominator. That choice
with \(s=1\) makes the denominator identically zero. The opposite signs
in the table follow the defining identity. For tangent and cotangent the
notebook also sets \(P=Q\) and \(c=f=-d=-g\). Either selector then uses the
same polynomial argument in both halves of the ratio.
\src{pp.~58--65}

\subsection{Where the derivation stops}
The final note identifies \(Y(x;u,V)\) with a component of a polynomial field
whose coefficient and exponent tables have been set to \(u,V\). Along a
solution of that polynomial system, time-differentiating its right-hand side
\(n-1\) times gives the \(n\)th derivative of the corresponding coordinate.
This observation explains the hope that the polynomial result could be
reused inside the exponentials.

There is a dependence on the trajectory in that last step. If the actual field
is \(\mathcal F\) from \eqref{eq:general}, then along its solution
\(\dd Y/\dd t=\sum_j(\partial_jY)\mathcal F_j\).
For example, \(\dd(e^C)/\dd t=e^C\sum_j(\partial_j C)\mathcal F_j\).
Differentiating a quotient further introduces products of these factors and
powers of its denominator. The argument derivative therefore uses the full
field governing the trajectory. Pages 58--65 establish the representation
\eqref{eq:general}. The notebook closes with the higher derivatives of the
transcendental contribution posed as a problem. Its completed numerical step
is \eqref{eq:step} for the autonomous polynomial system.
\src{pp.~59--65}

\section*{Machine checks and companion paper}
\addcontentsline{toc}{section}{Machine checks and companion paper}
The accompanying Lean~4 project gives a second proof of the polynomial
construction. Its variables range over arbitrary finite dimensions, column
counts and derivative orders. The theorem \file{Nomad.monomial_hessian}
checks the successive partial derivatives on notebook p.~2. With one
exponent vector per input column,
\file{Nomad.lie_table} and \file{Nomad.collect_table} give
\eqref{eq:rec} and \eqref{eq:group}. The theorem \file{Nomad.field_paths}
expands the same polynomial into the choices and products of \eqref{eq:path}.
A branch that differentiates a zero power has coefficient zero throughout
its descendants. The theorem \file{Nomad.slide_bounds} checks the overlap
limits for either ordering of the two table widths.

Given a solution of \eqref{eq:ivp} on a nondegenerate interval,
\file{Nomad.jet_derivative} identifies the polynomial at every order
with the corresponding time derivative, including the one-sided derivatives
at the endpoints. Theorems \file{Nomad.normalized_succ} and
\file{Nomad.step_eq_taylor} give the factorial scaling and the Taylor evaluation
in \eqref{eq:step}. In a separate file, \file{Nomad.selectors} verifies
the twelve choices in Table~\ref{tab:selectors} on their indicated domains.

The companion draft \emph{Local error bounds for consecutive NOMAD Taylor
steps} derives an error enclosure from the next two orders and gives a
PyTorch example with rational endpoints. Its theorem
\file{Nomad.boxed_consecutive} bounds the error using only the field and
derivative tables, a supplied box, and the step length, given a solution on the
step. To check all proofs, run \file{lake build} from the \file{lean/}
directory using its pinned Lean and mathlib v4.19.0 toolchain.

\begin{thebibliography}{9}
\addcontentsline{toc}{section}{References}
\raggedright
\bibitem{notebook}
J. Atkinson. \emph{NOMAD: N-th Order M-Dimensional Nonlinear Diff Eq Solver}.
Unpublished photographic notebook, \file{NOMAD_derivation.pdf}, 65 pages.
Clean review dated 27 August 2017.
\bibitem{jorba}
À. Jorba and M. Zou. A software package for the numerical integration of
ODEs by means of high-order Taylor methods. \emph{Experimental Mathematics}
\textbf{14} (2005), 99--117.
\href{https://doi.org/10.1080/10586458.2005.10128904}{doi:10.1080/10586458.2005.10128904}.
\bibitem{pal}
A. P\'al and \'A. S\"uli. Solving linearized equations of the \(N\)-body
problem using the Lie-integration method. \emph{Monthly Notices of the Royal
Astronomical Society} \textbf{381} (2007), 1515--1526.
\href{https://doi.org/10.1111/j.1365-2966.2007.12248.x}{doi:10.1111/j.1365-2966.2007.12248.x}.
\bibitem{carleman}
S. V. Chermnykh. Carleman linearization and normal forms for differential
systems with quasi-periodic coefficients. \emph{SpringerPlus} \textbf{5}
(2016), article 1347, Proposition 4.
\href{https://doi.org/10.1186/s40064-016-3015-6}{doi:10.1186/s40064-016-3015-6}.
\end{thebibliography}

\ifnotebooknotes
\appendix
\section{Notes in the order of the photographs}\label{app:leaves}
Page numbers below refer to the collated notebook \cite{notebook}. The main
text consolidates repeated algebra. This appendix keeps the working remarks,
changes of notation, discarded approaches and intermediate checks near their
original page numbers. Quotation marks denote short transcriptions. Other
remarks are rendered in the notation of this edition. The cover is p.~1.
The dated review begins on p.~3, immediately after the foundation sheet.

\notegroup{Foundation and review}
\leaf{2. The polynomial sheet.}{\(N\) is the order in a single variable and
\(M\) is the number of dimensions. The row of \(K\) identifies the variable,
the column its term. When differentiating in \(x_w\), retain precisely the
indices for which \(K_{wi}\ne0\). For the second derivative the cases \(y=w\)
and \(y\ne w\) have different coefficients and counts. The sheet draws the
gradient as a column and the Hessian as an \(M\times M\) matrix. See
\eqref{eq:grad}--\eqref{eq:hmixed}.}
\leaf{3. “Notation description.”}{The 27 August review constructs \(K\) in
blocks. In the original \(\hat k(d,b)\), each of \(N,N-1,\ldots,0\)
multiplies a row of \((N+1)^{d-1}\) ones, and the concatenated row is
repeated \(b\) times. Stack
\(\hat k(M,1),\hat k(M-1,N+1),\ldots,\hat k(1,(N+1)^{M-1})\).
The coefficient row selects which terms occur in that component.}
\leaf{4. “More Notation Description.”}{The review specifies an autonomous
multivariate polynomial system, differentiates it twice beyond the input
equation, and defines \(\delta_{ij}=\delta[i-j]\). The remark at the foot of
the page is that “the repeating summations are cumbersome to write!”}
\leaf{5. The multiple summation.}{There are \(n-1\) choices of \(j\) and
\(n-1\) of \(i\) in addition to \(\alpha\). In the dense table every \(i\)
has \((N+1)^M\) possibilities. At \(n=1\) both inner sums are marked
“(i.e. no summation).”}
\leaf{6. “Notation complete, solution.”}{The coefficients, the successive
powers brought down by differentiation, and the final power of each
\(x_\beta\) are grouped separately. A margin line spells out the third-order
exponent \(K_{\beta\alpha}+K_{\beta i_2}-\delta_{\beta j_2}
+K_{\beta i_1}-\delta_{\beta j_1}\). The last line is the general product
in \eqref{eq:three}.}
\leaf{7. “Series Solution.”}{Evaluate the derivatives at \(t_0\) and multiply
each by \((t-t_0)^n/n!\). The review estimates a dense box of
\([2^{n-1}N+1]^M\) terms and remarks that even small \(M\) makes this
allocation impractical at large \(n\). Many columns are absent in the input
because the coefficient table has “many zeros!”}
\leaf{8. Preparing the matrices.}{Parse the input field, choose its \(K\),
remove paired all-zero coefficient columns and call the survivors
\(\hat a(1),\hat k(1)\). The hats subsequently label the derivative order.
The contribution on the bottom line uses the old coefficient, the old power
and the \(j\)th field coefficient. The differentiated power's row is clarified
in \eqref{eq:rec}.}
\leaf{9. “Alg. Continued.”}{Evaluate the \(n\)th derivative from its new
coefficient and exponent tables. The computational shortcut is “an operation
similar to convolution.” The concern is that computing the new tables
“term by term is horrendously inefficient.” The review directs the reader to the other notebook
for an example and proposes dividing the coefficient by the current order
as the Taylor series is constructed. It ends, “This concludes review of
original NOMAD.”}

\notegroup{Working out the derivatives}
\leaf{10. A single scalar polynomial.}{The original derivation begins again
with \(P(x)\) and the block construction of \(K\). If \(a_i=0\), its term
vanishes. This observation later becomes a whole-column test for a vector
of coefficients.}
\leaf{11. The problem statement.}{“Things are about to get messy.”
Each equation \(\dot x_\ell\) gets its own row \(a^\ell_i\) over the shared
\(K\). The boxed task is to find the derivative of arbitrary order for any
output coordinate \(\ell\).}
\leaf{12. A single product.}{“Sufficient to examine how” the monomial
evolves under repeated differentiation. The first application peels off
\(x_1^{k_1}\), leaving \(x_1^{k_1}(x_2^{k_2}\cdots x_M^{k_M})'\)
and \(k_1x_1^{k_1-1}\dot x_1x_2^{k_2}\cdots x_M^{k_M}\).
The trial formula at the bottom retains an extra copy of the original product.}
\leaf{13. Closing the product rule.}{The accumulated outer terms are called
\(F\) temporarily while the final two factors are differentiated. The
“Whoops...” correction accounts for exactly one differentiated factor per
summand. Substituting the polynomial for each \(\dot x_j\) yields
\eqref{eq:one}.}
\leaf{14. Replacing the inner index.}{The unmatched product is crossed out.
The note “\(i_\alpha\) is confusing notation” prompts a rewrite using a
separate column index \(\beta\) for the inserted right-hand side.}
\leaf{15. Combining powers.}{Expand the inner polynomial against the old
monomial. The two powers of each variable add, with a single Kronecker
subtraction in the differentiated coordinate. This puts the second derivative
back in polynomial form.}
\leaf{16. What still depends on time?}{Braces label the old coefficients
and powers “time indep” and the product in \(x\) “time dep.” The latter
alone has to be differentiated for the third derivative.}
\leaf{17. A second product expansion.}{The substitution \(\beta=i_\alpha\)
relocates the previous field-column index. The page spells out how each of
the \(M\) combined powers, in turn, brings down its exponent and a factor
\(\dot x_j\).}
\leaf{18. The third derivative.}{The last state derivative is expanded as
another field sum. The boxed expression contains three original powers in
each final exponent and two Kronecker subtractions. Its new scalar multiplier
is the combined exponent from the previous differentiation.}
\leaf{19. A prediction.}{The side note uses a superscript on \(a\) and \(k\)
as a row index. Keeping the state component as a subscript avoids confusing
it with an exponent. The fourth-derivative pattern is written before the
calculation and subsequently marked “correct!”}
\leaf{20. Testing the prediction.}{Set
\(\hat k=k_\alpha+k_{i_1}+k_{i_2}-e_{j_1}-e_{j_2}\).
The old expression becomes \(x^{\hat k}\). Its product expansion has exactly
the form already worked out on pp.~12--14.}
\leaf{21. Recognizing the pattern.}{Replace \(\dot x_{j_3}\) by its
polynomial and restore the full expression for \(\hat k\). The coefficient
gains \(\hat k_{j_3}A_{j_3 i_3}\) and the exponent gains
\(k_{i_3}-e_{j_3}\). The handwritten celebration ends with “the whole thing!”}
\leaf{22. “Testing before going farther.”}{For \(\dot x=y^2-x\), mark
\(a_7=1,a_6=-1\). For \(\dot y=x^2-y\), mark \(a_3=1,a_8=-1\).
All other entries in each row are zero. These are the one-based indices
in the full nine-column table of Section~\ref{sec:example}.}
\leaf{23. Counting sums.}{One initial column and two new indices for each
subsequent differentiation give \(2n-1\) summations. The sheet distinguishes
the \(M\)-valued \(j\)'s from the \((N+1)^M\)-valued \(i\)'s. The clean
review fixes the lengths of the two lists at \(n-1\).}
\leaf{24. Renaming the indices.}{The old outer \(i\) becomes \(\alpha\) and
the product index \(j\) becomes \(\beta\). Successive coordinate indices
become \(j_1,j_2,j_3\), with matched field columns \(i_1,i_2,i_3\);
\(\ell\) stays fixed. All four derivatives are rewritten under this convention.}
\leaf{25. Array evaluation.}{The columns of the state array repeat the same
state, a “structure only for convenience of code!” Elementwise powers followed
by column products give \(\bar X_\alpha=x^{k_\alpha}\). The page frames
\(a^\ell_\alpha\bar X_\alpha\) as a single contribution and its sum over
\(\alpha\) as the component derivative.}
\leaf{26. Factoring the second derivative.}{Alongside the first four
derivatives the notes factor a term of \(x_\ell''\) into two field summands and
\(K_{j\alpha}/x_j\). The extra column subscript on the identical state copies
can be dropped, which frees the \(i_1\)-sum to give \(F_j\).
Equation~\eqref{eq:one} keeps all powers nonnegative wherever the
differentiated term survives.}
\leaf{27. Factoring the third derivative.}{Arrows match each \(a\) with
a monomial \(\bar X\). Where the selected coordinates are nonzero the
summand can be written as the three component terms
\((A_{\ell\alpha}x^{k_\alpha})(A_{j_1i_1}x^{k_{i_1}})
 (A_{j_2i_2}x^{k_{i_2}})\)
times \(K_{j_1\alpha}(u_1)_{j_2}/(x_{j_1}x_{j_2})\).
Equation~\eqref{eq:d3} gives its polynomial form on the whole state space.}
\leaf{28. “Notation Guide.”}{The equality of the state copies and the
distinction between one summand and the full derivative are both marked
“critical importance.” The summed column labels cover the width of the
exponent table. \(M\) counts state rows.}
\leaf{29. Prior factors under a brace.}{The third and fourth derivatives
are annotated to show which factors come from the preceding derivative.
Set \(j_0=\ell,\ i_0=\alpha\), and suppress the initial delta correction
using a discrete step. A side note notices the awkward product limit at
\(n=1\).}
\leaf{30. The boxed general result.}{An empty product contributes one.
The displayed product of coefficients, product of exponent sums, and
monomial with accumulated powers is reproduced by
\eqref{eq:path}--\eqref{eq:three}. The field in the lower box is its
first-order base case.}

\notegroup{Finding the iterative computation}
\leaf{31. Regrouping as a partial derivative.}{The original \(a^\ell_\alpha\)
has to sit inside the derivative of the summand. Collecting the inner field
sum gives \(\sum_{\alpha,j}\partial_j(A_{\ell\alpha}x^{k_\alpha})F_j\).
The page ends, “Need recursive construction of higher order.”}
\leaf{32. Absorbing two summations.}{The \(i_2\)-sum is \(F_{j_2}\).
For a fixed previous term, the power brought down by \(\partial_{j_2}\)
is already known. Thus the \(j_1,i_1\) choices can be absorbed into a new
column label. The new coefficient is
\(A_{\ell\alpha}A_{j_1i_1}K_{j_1\alpha}\), paired with
\(k_\alpha+k_{i_1}-e_{j_1}\).}
\leaf{33. “Only 2 sums!”}{With each old polynomial already stored, the
chain rule sums over its terms and the differentiated coordinate. The
component \(F_j\) still has its own column sum when expanded. The observation in the margin is “only 2 sums! w/ recursive computation!” A recursive
loop sketch is followed by the word “iterative” over “recursive” and a
new aim, the repeated construction of coefficient and exponent tables.}
\leaf{34. The dense allocation.}{Successive maximum powers are pencilled
in as \(2N,4N,8N,\ldots\). The prohibitive size of the corresponding dense
box motivates an “adaptive algorithm that minimizes the size of \(K\)
based on coefficients \(a\).” The count and a tighter bound appear in
Section~\ref{sec:slide}.}
\leaf{35. Parsing and compression.}{Parse \(a\) in the ordering of the
\((M,N)\) exponent table. The note says “parser needs some more work, good
for now.” A column with zero in every output row can be removed from both
tables. The next task is to relate \(\hat a,\hat k\) to the derivative order.}
\leaf{36. Checking the second derivative again.}{The product rule is
worked from the factors \(x_1^{k_1},x_2^{k_2},\ldots\) back to
\(\sum_{\alpha,j}A_{\ell\alpha}K_{j\alpha}x^{k_\alpha-e_j}\dot x_j\).
This confirms which old exponent is used when filling a coefficient column.}
\leaf{37. When do two terms coincide?}{Two triples can share a new column
precisely when \(k_\alpha+k_i-e_j=k_{\alpha'}+k_{i'}-e_{j'}\) in every
coordinate. The corresponding products \(A_{\ell\alpha}A_{ji}K_{j\alpha}\)
then add. This question leads directly to the lookup rule.}
\leaf{38. A discarded matrix product.}{The page tries \(AK^\mathsf{T}A\)
and crosses it out. Transposing \(K\) relocates an entry without changing
which variable and monomial it denotes. Ordinary matrix multiplication
contracts the column label before the new exponent has been determined.}
\leaf{39. Block strides.}{The drawn rows of \(K\) are constant in descending
blocks. With stride \(s=(N+1)^{M-j}\), the pencilled identity
\(K_{ji}=K_{j,i+\gamma s}+\gamma\) applies while the digit stays within the
same descending cycle. Crossing a cycle wraps the digit, as in \eqref{eq:digit}.}
\leaf{40. The upper corner.}{Even when both input columns are
\((N,\ldots,N)\), differentiating lowers one coordinate. The exponent with
\emph{all} coordinates \(2N\) is therefore absent from the second derivative.
This identifies a guaranteed zero slot in the full box.}
\leaf{41. The lower corner.}{Choosing the constant as the differentiated
monomial gives zero because every pulled-down exponent vanishes.
This conclusion concerns that contribution. Another pair can produce a
constant output. For example, \(\dot x=1+x\) gives \(x''=1+x\).}
\leaf{42. Tabulating the candidates.}{Rows are arranged by fixed outer
column \(\alpha\) and chosen coordinate \(j\), listing every inserted column
\(i\). In the last outer block the differentiated monomial is the constant
column, so its whole block contributes zero.}
\leaf{43. Opening up the products.}{The first blocks of p.~42 are written
coordinate by coordinate. The \(-1\) moves from the exponent of \(x_1\)
to that of \(x_2\), and eventually to \(x_M\), as the chosen \(j\) changes.}
\leaf{44. Looking up collisions.}{Monomials coincide as polynomials when
all exponents match. With the same outer column and \(j\ne g\), another
inserted column must satisfy \(k_f=k_i-e_j+e_g\).
On the full input grid this requires \(K_{ji}\ge1\) and
\(K_{gi}\le N-1\). Equality at one numerical state would be insufficient
for collecting polynomial coefficients.}
\leaf{45. A two-dimensional output grid.}{For \(M=N=2\), the two input
tables have nine columns and a dense second-derivative box has 25. The
upper neighbors \((4,3)\) and \((3,4)\) are marked “highest unique.”
The corner \((4,4)\) is unreachable. The tentative mark at \((0,0)\)
is subject to the distinction on p.~41.}
\leaf{46. Coordinates constant within a block.}{The sheet pursues a
column-index formula through repeated first-row values. For two variables,
columns \(u(N+1)+1,\ldots,(u+1)(N+1)\) all have first coordinate \(N-u\),
\(0\le u\le N\). Equation~\eqref{eq:digit} gives the general row stride.}
\leaf{47. Shifting a row.}{Two copies of the nine-column table are aligned.
Subtracting \(e_1\) from the field columns changes their first row to
\(1,1,1,0,0,0,-1,-1,-1\), leaving the second row untouched.
The subsequent addition must use the aligned current column.}
\leaf{48. The first slide.}{The sideways photograph introduces the circled
plus, the direction of travel, and “Lookup \(\to\) Store.” The next alignment
can produce two different output locations. The note calls for a manual
test with the reduced matrices of \eqref{eq:xy}. Its index lists are
typeset in Section~\ref{sec:slide}.}
\leaf{49. The \(j=1\) pass.}{The page traces all seven alignments of the
reduced four-column tables. It crosses out invalid exponents and discards
zero vectors. The four surviving keys and their coefficient vectors are
\((1,0):(1,0)\), \((0,2):(-1,0)\), \((1,2):(0,2)\), and
\((2,0):(0,-2)\).}
\leaf{50. Unequal widths.}{The next slide uses the six columns of the
second-derivative table against four field columns. The array called
\emph{ali} holds \(\alpha\) and \(i\) in two rows. The edges have increasing
and decreasing lengths, with three full overlaps in the middle. The same
sheet records the four-by-four case for comparison.}
\leaf{51. The \(j=2\) pass.}{Here the surviving keys and vectors are
\((2,0):(0,-1)\), \((2,1):(2,0)\), \((0,1):(0,1)\), and
\((0,2):(-2,0)\). Combine these with p.~49 and read off \eqref{eq:xyD2}.
The page marks the result “Correct!”}
\leaf{52. Counting and locating the tail.}{For widths \(P\ge Q\) the full
list has \(PQ\) columns. Before the decreasing overlaps begin, the cumulative
count is \(\sum_{r=1}^{Q-1}r+Q(P-Q+1)=Q(2P-Q+1)/2\).
The handwritten offset checks give 10 for \(P=Q=4\) and 18 for \(P=6,Q=4\).}
\leaf{53. The algorithm in words.}{Slide \(K-e_j{\bf1}^\mathsf{T}\) across
the current exponent table for every \(j\), keeping the pair of indices. “If the sum is valid” in all coordinates, compute its coefficient from the
indices. The instruction for equal exponents is to “sum” their coefficients
“up if it repeats.” Remove entries with zero vectors, then sort the remaining exponents
“while maintaining the pairs.”}

\notegroup{Scaling the series}
\leaf{54. Finite precision.}{“Numerical issues when number of derivatives
get large!” The notes speak of the spacing of large values representable
in double precision. Hats denote the new derivative table, tildes its
predecessor, and the plain \(a,k\) denote the input field. The explicit
path expression is inserted into the Taylor series at \(x=x(t_0)\).}
\leaf{55. The source of large numbers.}{Compare the compact derivative sum
with its expanded triple sum over \(\alpha,j,i\). Both the raw coefficients
and \(n!\) can become very large. The margin associates the concern with
orders around 50--100. This observation about scale depends on the field
and the arithmetic.}
\leaf{56. Unfolding the division.}{Expand \(n!=n(n-1)\cdots 2\) as the
coefficient construction is traced backward. The last line displays the
initial coefficient divided by \(1!\), followed by factors divided by
\(2,3,\ldots,n\). The two equivalent placements of the final divisor
are given in Section~\ref{sec:taylor}.}
\leaf{57. “The factorial will happen gradually.”}{The displayed series
uses a final \(1/n\) with a partially scaled table. The accompanying note
says that the division will “help control the growth of the elements” in the
coefficient table and keep them within useful floating-point range.}

\notegroup{The extension sheets}
\leaf{58. Regular and hyperbolic functions.}{Use exponentials of \(jx\) and
\(-jx\) for sine, cosine and tangent, then use imaginary arguments for sinh,
cosh and tanh. The notebook writes \(j=\sqrt{-1}\), rendered \(\ii\) in
Section~\ref{sec:trans} to free \(j\) for coordinate indices.}
\leaf{59. A new term.}{The first proposed term is
\(b[(e^{Y_1}+(-1)^re^{-Y_1})/(e^{Y_2}+(-1)^se^{-Y_2})]^h\).
\(Y_1,Y_2\) are multivariate polynomials. Bits \(r,s\) select a sum or a
difference in each pair. The note about weighted exponential combinations
uses a quotient for tangent and the reciprocal functions.}
\leaf{60. Four polynomial exponents.}{The sheet writes the coefficients
\(c,d\) over a shared table \(P\) and \(f,g\) over \(Q\), then substitutes
the full polynomial sums into the quotient. An early sine prescription is
crossed out. The surviving shorthand \(Y(x;u,V)\) counts the monomials by
the cardinality of the coefficient row.}
\leaf{61. Selecting sine.}{Set \(b=0\) to omit the transcendental contribution.
For sine take \(f=g=0,\ s=0\), hence denominator \(2\). Take \(r=1\),
opposite purely imaginary \(c,d\), and a purely imaginary prefactor. The
margin identifies the desired real amplitude as the coefficient of sine.
Table~\ref{tab:selectors} makes its sign explicit.}
\leaf{62. Selecting sinh, cos and cosh.}{All three keep the denominator at
two. Sinh uses a difference of exponentials with opposite real arguments.
Cosine and cosh use sums, with opposite imaginary and real arguments,
respectively. The prefactor is real. The sheet asks for at least one nonzero argument
coefficient in its selection rules.}
\leaf{63. Selecting reciprocals.}{Set \(c=d=0,\ r=0\) so the numerator
equals two. An even denominator yields sec or sech, and an odd denominator
yields csc or csch. Circular cosecant uses a purely imaginary prefactor.
The real-denominator sign for csch is the correction recorded after
Table~\ref{tab:selectors}.}
\leaf{64. Selecting tan, tanh, cot and coth.}{Match the argument polynomials
in numerator and denominator by taking \(c=f,\ c=-d,\ f=-g,\ P=Q\).
The pair \((r,s)=(1,0)\) gives the tangent family, and \((0,1)\) gives the
cotangent family. Pure imaginary arguments yield the circular functions,
real arguments the hyperbolic ones. The sheet concludes, “All typical forward regular \& hyperbolic trig functions
can be expressed w/ this new term!”}
\leaf{65. “General equation.”}{The top line includes both \(a\) and the
powered quotient in the coefficient of each monomial. The starred system
\(x^*\) denotes the polynomial case for which the \(n\)th derivative is
known. The lower box sets out the remaining transcendental part and
identifies \(Y(u,V)\) as a polynomial right-hand side at \(a=u,\ k=V\).
The arrows end at this identification. Section~\ref{sec:trans} states which
trajectory is required when it is differentiated.}

\paragraph{Editorial conventions.}
\(A_{\ell i},K_{\beta i}\) render the notebook's \(a_i^\ell,k_i^\beta\).
The induction proof, sharper degree bound, error discussion, independent
arithmetic check and literature comparisons are editorial additions.
Substantive readings concern the differentiated row (p.~8), constant
terms (pp.~41, 45), factorial scaling (pp.~56--57) and the csch sign (p.~63).

\clearpage
\section{Implementation notes}\label{app:code}
The MATLAB comparison uses the supplied sources. The PyTorch comparison records
the port as received and the subsequent corrections. Numerical checks used
PyTorch 2.6.0 on CPU. The author reports faster
execution with the PyTorch port. This edition supplies no comparative timings.

\paragraph{MATLAB.}
The recursive table constructor in \file{matlab/development_assistance_scripts/}
(\file{kExps.m}) generates the descending blocks of \eqref{eq:digit}.
The parser and the early Taylor routine are in
\file{matlab/backup_and_parser/}. In \file{nomadTaylor.m} the order-\(n\)
derivative is divided explicitly by \file{factorial(i)}.
In \file{matlab/optimized/}, \file{nomad_optim.m} constructs \emph{ali},
forms the shifted columns, groups equal rows of the transposed exponent
table, and reverses the \emph{columns} of both output matrices.
\file{nomadTaylor_optim.m} starts from the first derivative, divides the old
coefficient matrix by the new order when constructing orders \(2,\ldots,p\),
and evaluates every polynomial at the same symbolic initial state.

The zero-column test in \file{nomad_optim.m} takes the arithmetic sum
down the rows. A column such as \((1,-1)^\mathsf{T}\) then disappears even
though it contributes to both equations. The exact rule in the notes is
\(\forall\ell,\ c_\ell=0\), applied also after contributions with identical
exponents have been added. The three-region index construction in the code
assumes the current width is at least the field width, as on p.~50. The
offset limits in Section~\ref{sec:slide} also cover shrinking support.

\paragraph{PyTorch port as received.}
\file{pytorch/nomad_solver.py} precomputes the derivative tables, groups
exponents with \file{torch.unique}, and evaluates a batch of initial states.
Three operations in the historical port explain the second-order discrepancies.
At zero-based loop index \(i\), the constructor divided the previous
table by \(i\) while constructing derivative \(i+1\). This gives the partial
factorial convention on p.~57. The evaluation omitted the remaining factor.
After grouping, \file{torch.flip(..., dims=(0,))} reversed the coordinate
rows of both matrices. For \(F=(x^2,0)\), the raw second derivative
\((2x^3,0)\) became \((0,2y^3)\). Finally, \file{step} evaluated each term
at the partial sum from the preceding order. For \(F(x)=x\), \(z=1\), this
gave \(1+h+(1+h)h^2\), or \(1.328125\) at \(h=1/4\).
Working arrays also inherited the default dtype after the input tables had been
converted to double.

\paragraph{PyTorch corrections accompanying the formalization.}
The port now divides by the new order \(i+1\), reverses columns after
grouping, and evaluates every term at the beginning-of-step state.
The scalar example returns \(1+h+h^2/2=1.28125\) at \(h=1/4\).
For \(F=(x^2,0)\), the second Taylor coefficient is \((x^3,0)\).
The overlap limits in Section~\ref{sec:slide} handle shrinking supports, and
the coefficient buffers retain double precision. A positive pruning threshold
discards candidates before collection and changes the recurrence.
The comparison in \file{pytorch/test_nomad.py} measures mean trajectory error
on \eqref{eq:xy}. The new regressions check coefficients and steps directly,
with zero pruning threshold, in \file{pytorch/test_nomad_error_bounds.py}.

\fi
\end{document}
